% =====================================================================
% 03b_infrastructure_adhoc.tex
% Mo hinh Infrastructure va Ad-hoc (IBSS): doi chieu chi tiet.
% Nguon: "Mo hinh infrastructure.docx" (My Drive > Networking) va bo anh
% "picture source" (Picture1-Picture6). Duoc \input ngay sau
% sections/03_architecture.tex nen cac tieu de duoi day tiep tuc danh so
% trong cung Muc "Kien truc va cac thanh phan WLAN".
% =====================================================================
\begingroup
\newcolumntype{L}[1]{>{\hsize=#1\hsize\raggedright\arraybackslash}X}

\subsection{Mô hình hạ tầng (Infrastructure) và mô hình tùy biến (Ad-hoc)}

Trong phạm vi gia đình và doanh nghiệp nhỏ, mạng cục bộ không dây thường được
thiết lập theo hai mô hình chính: mô hình hạ tầng (Infrastructure, có thiết bị
trung tâm) và mô hình tùy biến hay ngang hàng (Ad-hoc). Mô hình hạ tầng là lựa
chọn phổ biến gần như mặc định, trong khi mô hình Ad-hoc trên các thiết bị đầu
cuối thông dụng đã dần được thay thế bởi các công nghệ ngang hàng mới hơn như
Wi-Fi Direct, và chỉ còn tồn tại bền vững trong những dạng tiến hóa chuyên biệt
(MANET, VANET, FANET) ở các lĩnh vực không có hạ tầng cố định.

\subsubsection{Mô hình hạ tầng (Infrastructure)}

Đây là một hệ thống phân cấp: một nhóm thiết bị đóng vai trò quản lý và cung cấp
dịch vụ, nhóm còn lại đóng vai trò thụ hưởng. Ba thành phần thiết bị chính gồm:

\begin{itemize}[leftmargin=*]
  \item \textbf{\acf{ap} / bộ định tuyến không dây:} phát sóng Wi-Fi (công bố
    \ac{ssid}), nhận và chuyển tiếp dữ liệu giữa các thiết bị không dây với nhau
    và giữa thiết bị không dây với mạng Internet nếu có. Trong hộ gia đình, thiết
    bị này thường tích hợp thêm chức năng modem (giải điều chế tín hiệu của nhà
    cung cấp dịch vụ) và switch (chia cổng mạng có dây).
  \item \textbf{Thiết bị đầu cuối (\acp{sta}):} mọi thiết bị có card mạng không
    dây (WNIC -- Wireless Network Interface Card) như máy tính xách tay, điện
    thoại thông minh, Smart TV, camera Wi-Fi hay bóng đèn thông minh; các trạm
    này chỉ sử dụng dịch vụ do \ac{ap} cung cấp.
  \item \textbf{Hạ tầng dây dẫn:} cáp mạng Ethernet RJ45 hoặc cáp quang nối bộ
    định tuyến ra bên ngoài, hoặc nối tới các switch khác trong tòa nhà; đây
    chính là \ac{ds} đã trình bày ở Mục~\ref{sec:arch}.
\end{itemize}

\textbf{Phạm vi định danh: \ac{bss} và \ac{bssid}.} \ac{bss} là đơn vị cấu
thành cơ bản nhất của mạng Wi-Fi theo mô hình hạ tầng: một \ac{ap} duy nhất cùng
toàn bộ thiết bị không dây đang liên kết trực tiếp với \ac{ap} đó. Để phân biệt
các \ac{bss} với nhau, mỗi \ac{bss} mang một mã định danh duy nhất là
\ac{bssid}, và trong mô hình hạ tầng \ac{bssid} chính là địa chỉ MAC vật lý của
giao diện vô tuyến trên \ac{ap}.

\textbf{Mở rộng và chuyển vùng: \ac{ess}.} \ac{ess} là mô hình mở rộng, hình
thành khi nhiều \ac{bss} (nhiều \ac{ap}) được nối với nhau qua \ac{ds} để phủ
sóng một không gian lớn hơn. Mọi \ac{ap} trong cùng một \ac{ess} dùng chung một
\ac{ssid} và cùng cấu hình bảo mật, mặc dù mỗi \ac{ap} vẫn sở hữu
\ac{bssid} riêng. Nhờ đó thiết bị đầu cuối có thể di chuyển liên tục giữa các
vùng phủ sóng mà không bị ngắt kết nối (roaming). Cần lưu ý rằng \ac{ess}
truyền thống không có cơ chế chuyển vùng nhanh, nên việc chọn \ac{ap} nào để
liên kết là do chính thiết bị của người dùng quyết định dựa trên cường độ tín
hiệu; các tiêu chuẩn 802.11k/v/r sau này mới bổ sung cơ chế hỗ trợ chuyển vùng
nhanh và cân bằng tải.

\subsubsection{Mô hình tùy biến (Ad-hoc / IBSS)}

Mô hình Ad-hoc là một mạng lưới phi tập trung: không có thiết bị nào làm chủ và
tất cả các nút đều bình đẳng. Thành phần thiết bị chỉ gồm các nút không dây
(\emph{wireless nodes}) -- chính là các thiết bị đầu cuối như máy tính xách tay
hay máy tính để bàn gắn card/USB Wi-Fi -- và tuyệt đối không có sự xuất hiện
của \ac{ap} hay bộ định tuyến. Mỗi nút vừa đóng vai trò máy khách (gửi yêu cầu)
vừa đóng vai trò máy chủ/bộ định tuyến (xử lý và chuyển tiếp gói tin cho nút bên
cạnh).

\textbf{Định danh trong \ac{ibss}.} Ở mạng hạ tầng (\ac{bss}), \ac{bssid} là địa
chỉ MAC thật của \ac{ap}. Ở mạng Ad-hoc, vì không có \ac{ap}, nút đầu tiên khởi
tạo mạng sẽ tự sinh một địa chỉ MAC ảo ngẫu nhiên để làm \ac{bssid}; các nút
tham gia sau đó cùng dùng chung mã \ac{bssid} ảo này để nhận diện nhau.

\textbf{Duy trì mạng bằng beacon luân phiên.} Trong \ac{bss}, \ac{ap} liên tục
phát khung beacon để thông báo mạng còn tồn tại. Trong \ac{ibss}, do các nút
bình đẳng, tất cả các nút tham gia sẽ luân phiên nhau phát beacon theo cơ chế
phân chia đều (beaconing chia đều), nhờ đó không cần một nguồn phát duy nhất mà
mạng vẫn được duy trì.

\textbf{Tính cô lập.} \ac{ibss} là một ``hòn đảo'' độc lập: mặc định nó không
thể giao tiếp với hệ thống mạng có dây bên ngoài và không có kết nối Internet.
Đây là khác biệt bản chất so với mô hình hạ tầng, nơi \ac{ds} và bộ định tuyến
biên mở đường ra mạng ngoài.

\textbf{Giới hạn công nghệ.} Mạng Ad-hoc hoạt động ổn định nhất trên các chuẩn
Wi-Fi cũ là 802.11b, 802.11g và 802.11n ở băng tần 2,4~GHz. Về bảo mật, phần
lớn triển khai \ac{ibss} trên hệ điều hành phổ thông chỉ hỗ trợ mã hóa
\ac{wep} -- loại mã hóa đã bị chứng minh có thể bị phá trong vài phút -- hoặc
buộc phải để mạng mở không mật khẩu; do đó mô hình này không phù hợp cho dữ liệu
nhạy cảm. Ngoài ra, hiệu năng suy giảm rõ rệt khi số nút tăng, vì toàn bộ việc
điều phối kênh diễn ra phân tán và không có thiết bị nào quản lý tài nguyên vô
tuyến.

\textbf{Tiến hóa của mô hình.} Trên các thiết bị hằng ngày như điện thoại, máy
tính xách tay hay máy tính cá nhân, chế độ Ad-hoc đã được thay thế bằng
Wi-Fi Direct. Đối với các mạng như MANET, VANET hay FANET, đây là phiên bản
tiến hóa cao cấp nhất của mạng Ad-hoc: chúng vẫn thuần túy ngang hàng nhưng được
bổ sung các giao thức định tuyến (routing protocol) rất phức tạp ở tầng mạng
(lớp~3) để chuyển tiếp gói qua nhiều chặng, và vẫn đang được sử dụng trong những
lĩnh vực chuyên biệt như cứu hộ -- cứu nạn, giao thông thông minh hay thiết bị
bay không người lái.

\subsubsection{Đối chiếu chi tiết giữa hai mô hình}

Bảng~\ref{tab:infra-adhoc-detail} hệ thống hóa các khác biệt ở mức kiến trúc,
định danh, vận hành và bảo mật. Bảng này bổ sung cho bảng so sánh tổng quát
(Bảng~\ref{tab:infra-vs-adhoc}) bằng các tiêu chí cụ thể đã được kiểm chứng
trong quá trình thực nghiệm.

\begin{table}[htbp]
  \centering\footnotesize
  \caption{Đối chiếu chi tiết giữa mô hình hạ tầng (Infrastructure) và mô hình
  tùy biến (Ad-hoc / IBSS).}
  \label{tab:infra-adhoc-detail}
  \begin{tabularx}{\linewidth}{@{} L{0.75} L{1.15} L{1.1} @{}}
    \toprule
    \textbf{Tiêu chí} & \textbf{Mô hình hạ tầng (Infrastructure)} & \textbf{Mô hình tùy biến (Ad-hoc / \ac{ibss})} \\
    \midrule
    Thiết bị trung tâm & Bắt buộc có \ac{ap} hoặc bộ định tuyến không dây đảm
      nhiệm phát sóng và chuyển tiếp & Không có; các \ac{sta} hoàn toàn ngang
      hàng và tự chuyển tiếp cho nhau \\
    Đơn vị cấu thành & \ac{bss} (một \ac{ap}) mang \ac{bssid} riêng; nhiều
      \ac{bss} ghép thành \ac{ess} & \ac{ibss}: một nhóm nút độc lập, không có
      \ac{ap} \\
    Nguồn \ac{bssid} & Địa chỉ MAC vật lý của giao diện vô tuyến trên \ac{ap} &
      Địa chỉ MAC ảo sinh ngẫu nhiên bởi nút khởi tạo mạng, các nút khác dùng
      chung \\
    Phát beacon & \ac{ap} phát định kỳ, một nguồn phát duy nhất & Các nút luân
      phiên phát beacon (beaconing chia đều) \\
    Điều phối truy cập kênh & \ac{ap} giữ vai trò chủ; có thể hỗ trợ DCF hoặc
      PCF/HCCA & Phân tán hoàn toàn, chỉ hoạt động theo DCF (\ac{csmaca}) \\
    Cấp phát địa chỉ IP & Máy chủ DHCP tập trung trên \ac{ds} & Gán tĩnh thủ
      công hoặc tự cấp địa chỉ nội bộ (APIPA / link-local) \\
    Kết nối ra mạng ngoài & Có, qua \ac{ds} và bộ định tuyến biên & Mặc định
      không có; mạng bị cô lập với \ac{lan} có dây và Internet \\
    Mở rộng và chuyển vùng & Tốt; thêm \ac{ap} cùng \ac{ssid} là mở rộng
      \ac{ess}, hỗ trợ chuyển vùng & Kém; mọi nút phải trùng \ac{ssid} và trùng
      kênh, không có chuyển vùng; hiệu năng giảm khi số nút tăng \\
    Bảo mật & Đầy đủ; hỗ trợ WPA2/WPA3, 802.1X/\ac{radius} cho doanh nghiệp &
      Hạn chế; thực tế chỉ \ac{wep} hoặc mạng mở, không có WPA2/WPA3 \\
    Chuẩn và băng tần & Mọi thế hệ 802.11 và mọi băng tần được cấp phép & Ổn
      định nhất trên 802.11b/g/n, băng tần 2,4~GHz \\
    Hỗ trợ trên hệ điều hành phổ thông & Mặc định trên Windows, Android, iOS và
      macOS & Windows 10/11, Android và macOS không còn hỗ trợ; chỉ Linux hỗ trợ
      qua nl80211/iw \\
    Ứng dụng thực tế & Gia đình, văn phòng, trường học, doanh nghiệp & Chia sẻ
      dữ liệu tạm thời, cứu nạn khẩn cấp, mạng cảm biến; nay chủ yếu được thay
      bằng Wi-Fi Direct \\
    \bottomrule
  \end{tabularx}
\end{table}

Từ đối chiếu trên có thể rút ra định hướng lựa chọn mô hình: hiệu năng, khả năng
mở rộng, kết nối Internet và bảo mật hiện đại là những lý do khiến mô hình hạ
tầng vẫn là mặc định cho mọi nhu cầu khai thác thực tế; mô hình Ad-hoc chỉ còn
giá trị ở những tình huống không thể dựng hạ tầng, hoặc trong các dạng tiến hóa
đa chặng phục vụ mục đích chuyên biệt. Chính vì giới hạn về bảo mật và khả năng
kết nối, phần thực nghiệm của báo cáo chỉ dùng chế độ \ac{ibss} để kiểm chứng
nguyên lý ngang hàng ở mức hai nút, và thực hiện trên nền Linux -- hệ điều hành
duy nhất trong nhóm phổ thông còn hỗ trợ đầy đủ chế độ này
(Mục~\ref{sec:adhoc-config} và Mục~\ref{sec:experiment}).

\endgroup
